% !TEX program = pdflatex
% Modelo visual: Apresentacao 15102026 - dados preliminares.
% Fontes e validacao: 2_INTI/audit/apresentacao_27102026/RELATORIO.md.
% Resultados existentes: testes_chap1 (S01, M01, M02 e comparacao M01C).
% Sintese complementar: testes2, testes4, testes5 e testes7.
\documentclass[aspectratio=169,11pt]{beamer}
\usetheme{Madrid}
\usepackage[T1]{fontenc}
\usepackage[utf8]{inputenc}
\usepackage[brazilian]{babel}
\usepackage{lmodern}
\usepackage{booktabs,array,tabularx}
\usepackage{amsmath}
\definecolor{verde}{HTML}{24574D}
\definecolor{azul}{HTML}{1769A5}
\definecolor{grafite}{HTML}{253A40}
\setbeamercolor{structure}{fg=verde}
\setbeamercolor{title}{bg=verde,fg=white}
\setbeamercolor{frametitle}{bg=verde,fg=white}
\setbeamercolor{normal text}{fg=grafite}
\setbeamertemplate{navigation symbols}{}
\setbeamertemplate{footline}{\leavevmode\hbox{\begin{beamercolorbox}[wd=.86\paperwidth,ht=2.7ex,dp=1.2ex,leftskip=1em]{author in head/foot}Retornos, instituições e conversão da vegetação nativa\end{beamercolorbox}\begin{beamercolorbox}[wd=.14\paperwidth,ht=2.7ex,dp=1.2ex,center]{date in head/foot}\insertframenumber\,/\,\inserttotalframenumber\end{beamercolorbox}}}
\setbeamersize{text margin left=8mm,text margin right=8mm}
\setbeamertemplate{itemize items}[circle]
\newcommand{\fonte}[1]{\par\vfill{\fontsize{7.5}{9}\selectfont\color{gray}Fonte: #1\par}}
\newcommand{\tpar}{\par\vspace{2mm}}
\newcolumntype{P}[1]{>{\raggedright\arraybackslash}p{#1}}
\newcolumntype{L}{>{\raggedright\arraybackslash}X}
\newcommand{\est}[2]{\ensuremath{#1}\;{\color{gray}(#2)}}
\renewcommand{\arraystretch}{1.15}
\title[Retornos, instituições e conversão]{Retornos Econômicos, Instituições e Conversão da Vegetação Nativa em Mato Grosso}
\subtitle{Regressões preliminares: OLS e limites da inferência}
\author{Fernando Souza de Vieira}
\institute{Tópicos em Microeconomia}
\date{27 de outubro de 2026}
\begin{document}

% ================================================
% #1
\begin{frame}
\titlepage
\end{frame}
%% COMENTÁRIO: 
% (1) Bom, hoje eu vou mostrar as primeiras regressões do meu trabalho sobre a conversão da vegetação nativa em Mato Grosso.
% (2) A ideia é investigar como os incentivos econômicos, a capacidade produtiva da terra e a fiscalização se relacionam com a abertura de áreas para agricultura e pastagem.
% (3) Vou apresentar a hipótese, explicar as variáveis e comparar os resultados, discutindo o que a OLS permite dizer nesta etapa e o que ainda precisa ser investigado.

% ================================================
% #2 -- H1 vigente: 0_introduction.tex, REF7. Funcao conceitual ja existente.
\begin{frame}{Hipótese e função}
\begin{block}{H1: retornos, instituições e conversão}
Elevações nos retornos econômicos esperados dos usos agropecuários aumentam a taxa de conversão da vegetação nativa, sendo essa resposta, \textit{ceteris paribus}, amplificada pela capacidade produtiva local da terra e modulada pela efetividade do ambiente institucional.
\end{block}
\medskip
\[
\begin{aligned}
\text{Conversão}=f\bigl(&\text{preços agropecuários},\\
&\text{produtividade da terra},\\ &\text{ambiente institucional}\bigr)
\end{aligned}
\]
% \small
% A função organiza o mecanismo proposto. Nos testes, preços, aptidão potencial e autuações representam dimensões observáveis desse mecanismo.
% \tpar
% \textbf{Alcance preliminar:} os três modelos OLS são aditivos e ainda não testam a amplificação ou a modulação previstas em H1.
% \fonte{Capítulo 1, Introdução, REF7--REF9; apresentação de 15/10/2026.}
\end{frame}
%% COMENTÁRIO:
% (1) A hipótese que eu quero testar é que, quando aumenta o retorno esperado da atividade agropecuária, cresce o incentivo para abrir novas áreas e, com isso, a conversão de vegetação nativa.
% (2) Espero que essa resposta seja maior onde a terra tem mais capacidade de produzir. Já as regras ambientais e a forma como são aplicadas podem mudar esse incentivo, porque entram no custo e no risco de abrir uma nova área.
% (3) Essa função resume as três dimensões que eu vou observar: preços, capacidade produtiva da terra e ambiente institucional. Os preços ajudam a representar os incentivos econômicos, embora não mostrem, sozinhos, quanto o produtor espera ganhar.
% (4) Nessas primeiras regressões, eu começo olhando as associações de cada variável com a conversão. Para ver se a associação entre preços e conversão muda conforme a aptidão ou a fiscalização, eu preciso olhar os modelos com interações.

% ================================================
% #3 -- LEIA_ME.md, auditoria_cobertura.csv e modelos_resumo.    csv.
\begin{frame}{Y: Qual parcela da vegetação nativa foi convertida?}
Uma taxa por município e ano, calculada separadamente para \textbf{agricultura} e \textbf{pastagem}.
\medskip
\[
\text{Taxa de conversão (\%)}=
\frac{\text{hectares nativos convertidos}}{\text{hectares nativos no início do ano}}\times100
\]
\par\medskip
\textbf{Exemplo ilustrativo:} 10 ha convertidos de 1.000 ha nativos = \textbf{1\%}.
\par\bigskip
\small
\textbf{Amostras usadas nas regressões}\par\smallskip
\begin{tabularx}{\textwidth}{L r r l}
\toprule
\textbf{Modelos} & \textbf{Municípios} & \textbf{Município--ano} & \textbf{Período}\\
\midrule
S01 (OLS simples) e\newline M01 (OLS múltipla) & 137 & 2.772 & 2001--2024\\
M02 (OLS com fiscalização) & 128 & 1.140 & 2002--2024\\
\bottomrule
\end{tabularx}
\par\medskip
\textbf{Painel desbalanceado:} entram os anos com dados completos nas variáveis do modelo. Cada município pode contribuir com uma quantidade diferente de anos.
% \fonte{MapBiomas 10.1; testes\_chap1, LEIA\_ME e modelos\_resumo.}
\end{frame}
%% COMENTÁRIO:
% (1) Aqui eu defino o que estou tentando explicar: a parcela da vegetação nativa de cada município que foi convertida em um ano. Faço essa conta separadamente para agricultura e para pastagem.
% (2) Eu divido a área convertida pela vegetação nativa que existia no início do ano e multiplico por cem. No exemplo, converter 10 dos 1.000 hectares dá uma taxa de 1%.
% (3) Cada observação é um município em um determinado ano. Nas regressões simples e múltipla, eu uso 2.772 observações de 137 municípios. Com a fiscalização, são 1.140 observações de 128 municípios, porque também preciso da informação válida do ano anterior.
% (4) Eu aproveito os anos em que as variáveis daquele modelo estão completas, sem exigir que o município tenha dados para o período inteiro. Por isso, cada município pode contribuir com uma quantidade diferente de anos: é o que eu chamo aqui de painel desbalanceado.

% ================================================
% #4 -- utilitarios.py, MAPA/ler; escalas.csv.
\begin{frame}{X: Variáveis explicativas}
\small
\begin{tabularx}{\textwidth}{P{1.7cm} L}
\toprule
\textbf{Símbolo} & \textbf{Medida utilizada}\\
\midrule
$P^k_{t-1}$ & Preços apresentados: soja ($s$) e milho ($m$). $P$ indica o preço do produto $k$ no ano anterior ($t-1$), em dólares de 2024 por tonelada.\\[1.5mm]
$A_i^c$ & Aptidão ($A$) ou potencial capacidade produtiva em cada cidade ($i$), medida pelo rendimento potencial da cultura ($c$) para soja e milho, em kg/ha.\\[1.5mm]
$F_{i,t-1}$ & Autos/Multas ($F$) do IBAMA em cada cidade ($i$) no ano anterior ($t-1$), divididos por 10.\\[1.5mm]
$B_i$ & Indicadores do bioma ($B$) -- Amazônia (referência), Cerrado ou Pantanal -- predominante no município ($i$).\\
\bottomrule
\end{tabularx}
\smallskip
\textbf{Padronização:} apenas preços e aptidões. $z(X)=(X-\bar X)/\sigma_X$ permite comparar coeficientes por \textbf{1 desvio-padrão}, apesar das unidades diferentes.
\par\smallskip
Médias e desvios usam \textbf{24 anos} (preços) e \textbf{141 municípios} (aptidões). Ficam fixos entre amostras para \textbf{manter a mesma unidade de comparação entre modelos}.
\tpar
Preços e aptidão não medem o retorno líquido esperado. Autuações não medem diretamente a efetividade institucional.
% \fonte{Banco Mundial, GAEZ e IBAMA; testes\_chap1, utilitarios.py e escalas.csv.}
\end{frame}
%% COMENTÁRIO:
% (1) Aqui eu mostro as variáveis que entram para explicar a conversão. Os preços apresentados são os de soja e milho no ano anterior, em dólares de 2024 por tonelada. O k só identifica qual produto eu estou olhando.
% (2) Para representar a aptidão da terra, eu uso o rendimento potencial do GAEZ, uma base da FAO e do IIASA. É uma estimativa de quanto a terra poderia produzir de soja ou milho por hectare, em peso seco dos grãos.
% (3) A fiscalização entra pela contagem de autos de infração de flora do IBAMA no ano anterior. Eu divido por dez para ler o coeficiente a cada dez autos adicionais. Essa divisão só muda a unidade de medida.
% (4) O B representa o bioma predominante do município. Eu incluo indicadores para Cerrado e Pantanal, usando a Amazônia como referência, para considerar diferenças médias de conversão entre os biomas.
% (5) Só preços e aptidão foram padronizados. Assim, eu leio seus coeficientes para uma variação de um desvio-padrão. Uso os 24 anos para preços e os 141 municípios para aptidão, mantendo médias e desvios fixos entre amostras para comparar os modelos na mesma escala.
% (6) O cuidado é que preço e aptidão, sozinhos, não mostram o retorno líquido que o produtor espera ter. Da mesma forma, contar autuações não me diz, por si só, quanto a fiscalização consegue evitar de conversão.

% ================================================
% #5 -- Formulas literais dos notebooks ols_simples e ols_multiplo.
\begin{frame}{O que entra em cada modelo}
\small
\textbf{Variável dependente:} taxa de conversão da vegetação nativa.
\par\medskip
\begin{tabularx}{\textwidth}{P{2.8cm} L P{2.5cm}}
\toprule
\textbf{Modelo} & \textbf{Variáveis incluídas} & \textbf{Destino apresentado}\\
\midrule
\textbf{S01} (simples) & Preço da soja & Agricultura\\[2mm]
\textbf{M01} (múltipla) & Preços de soja e milho + aptidões + bioma & Agricultura e pastagem\\[2mm]
\textbf{M02} (fiscalização) & \textbf{M01 + autos de flora do IBAMA} no ano anterior & Agricultura e pastagem\\
\bottomrule
\end{tabularx}
\par\smallskip
Agricultura e pastagem são estimadas separadamente.
\tpar
{\footnotesize M01/M02 mantêm o controle adicional de preço da especificação original. As tabelas exibem coeficientes selecionados, \textbf{sem reestimação}.\par}
\tpar
\textbf{Em comum:} constante, sem efeitos fixos municipais/anuais e sem interações.
\tpar
\textbf{Erros-padrão:} agrupados por município e ano. Sem correção por múltiplos testes.
% Equações originais preservadas como referência; nenhuma especificação foi alterada.
% L^d_{it}=\alpha_d+\sum_{k\in\{s,m\}}\beta_{kd}z(P^k_{t-1})+\sum_{c\in\{s,m\}}\theta_{cd}z(A_i^c)+\delta_d'B_i+\lambda_d Z_{t-1}.
% S01: y^{\mathrm{ag}}_{it}=\alpha+\beta_s z(P^s_{t-1})+u^{\mathrm{ag}}_{it}.
% M01: y^d_{it}=L^d_{it}+u^d_{it}.
% M02: y^d_{it}=L^d_{it}+\gamma_d F_{i,t-1}+u^d_{it}.
% Z_{t-1}: controle adicional de preço padronizado, mantido na especificação original e não exibido nas tabelas.
% d indica agricultura ou pastagem; u^d_{it} reúne os fatores não observados.
% Inferência: referência t com 23 ou 22 graus de liberdade.
% \fonte{testes\_chap1, ols\_simples.ipynb, ols\_multiplo.ipynb e utilitarios.py.}
\end{frame}
%% COMENTÁRIO:
% (1) A lógica aqui é começar com uma relação simples e ir acrescentando variáveis. No S01 apresentado, eu relaciono o preço da soja à conversão para agricultura.
% (2) No M01, eu incluo preços, aptidões e bioma, preservando o controle adicional de preço da estimação original. No M02, acrescento os autos de flora do IBAMA do ano anterior.
% (3) Para M01 e M02, eu faço uma regressão para agricultura e outra para pastagem. O que muda é a taxa que estou tentando explicar. Dentro de cada especificação, as variáveis explicativas são as mesmas para os dois destinos, mas os coeficientes são estimados separadamente.
% (4) Ao incluir fiscalização, também diminui a amostra. Então, uma mudança no coeficiente entre M01 e M02 pode vir tanto da nova variável quanto da mudança nas observações que entram na regressão. Vou separar essas duas coisas mais adiante.
% (5) Por enquanto, eu uso modelos sem efeitos fixos municipais ou anuais e sem interações. Os erros-padrão consideram a dependência dentro dos municípios e dos anos. A leitura dos coeficientes continua sendo de associação.

% ================================================
% #6 -- Fonte numerica: coef_S01_agricultura, coef_M01_agricultura, coef_M02_agricultura.
\begin{frame}{OLS: conversão para agricultura}
\small Coeficientes (erros-padrão), em pp. Preços e aptidões em desvios-padrão.
\par\smallskip
\begingroup
\fontsize{9.5}{11.4}\selectfont
\renewcommand{\arraystretch}{1.00}
\setlength{\tabcolsep}{4pt}
\begin{tabularx}{\textwidth}{L c c c}
\toprule
 & \textbf{S01: simples} & \textbf{M01: múltipla} & \textbf{M02: fiscalização}\\
\midrule
Preço da soja & \est{-0{,}0512}{0,0251} & \est{-0{,}0674}{0,0667} & \est{-0{,}0821}{0,0773}\\
Preço do milho & -- & \est{0{,}0481}{0,0602} & \est{0{,}0703}{0,0664}\\
Aptidão da soja & -- & \est{-0{,}0367}{0,0222} & \est{-0{,}0197}{0,0137}\\
Aptidão do milho & -- & \est{-0{,}0241}{0,0170} & \est{-0{,}0096}{0,0077}\\
Autos de flora / 10 & -- & -- & \est{-0{,}0040}{0,0058}\\
\midrule
Controles de bioma & Não & Sim & Sim\\
Observações & 2.772 & 2.772 & 1.140\\
Municípios / anos & 137 / 24 & 137 / 24 & 128 / 23\\
$R^2$ & 0,0185 & 0,0562 & 0,0404\\
\bottomrule
\end{tabularx}
\endgroup
\par\smallskip
\footnotesize Constante nos três modelos. M01/M02 mantêm bioma e controle adicional de preço, cujos coeficientes não são exibidos.
\par\smallskip\small
Nenhum preço ou aptidão rejeita zero a 5\%. Em S01, soja tem $p=0{,}053$: o sinal negativo não permite concluir que preços reduzam a conversão.
% \fonte{testes\_chap1, modelos\_coeficientes.csv e modelos\_resumo.csv.}
\end{frame}
%% COMENTÁRIO:
% (1) Aqui eu mostro os resultados para a conversão em agricultura. Os coeficientes estão em pontos percentuais, e os erros-padrão, entre parênteses, indicam a incerteza. Vou traduzir os valores mantendo as demais variáveis constantes. Os preços estão em dólares de 2024.
% (2) No S01, a associação com o preço da soja é negativa, com p igual a 0,053, próximo do corte de 5%. Isso não torna o coeficiente inválido, mas a interpretação continua sendo de associação, sem afirmar que o preço causou uma redução da conversão.
% (3) No M01, para cada US$ 132,44 por tonelada a mais no preço da soja do ano anterior, a taxa estimada de conversão é 0,0674 ponto percentual menor. Para cada US$ 73,57 por tonelada a mais no preço do milho do ano anterior, a taxa estimada é 0,0481 ponto percentual maior.
% (4) Para a aptidão, eu comparo municípios. Um município com rendimento potencial de soja 379,58 kg por hectare maior tem uma taxa estimada 0,0367 ponto percentual menor. Para o milho, uma aptidão 778,52 kg por hectare maior está associada a uma taxa 0,0241 ponto percentual menor. Esses rendimentos são medidos em peso seco dos grãos.
% (5) Essas diferenças de preço e aptidão equivalem a um desvio-padrão. Num exemplo hipotético de M01, se a taxa prevista fosse de 1%, um preço da soja US$ 132,44 por tonelada maior corresponderia a uma taxa prevista de aproximadamente 0,9326%.
% (6) No M02, cada dez autos de flora adicionais no ano anterior estão associados a uma taxa estimada 0,0040 ponto percentual menor. Essa estimativa também é incerta e não permite concluir se a fiscalização foi eficaz.
% (7) Os intervalos de confiança de 95% dos coeficientes que eu comentei incluem zero: as magnitudes são incertas. Mesmo um p-valor menor não bastaria para estabelecer causalidade nestes modelos. O R quadrado fica entre cerca de 2% e 6%: os modelos explicam pouca variação da conversão, mas isso, sozinho, não invalida a OLS.

% ================================================
% #7 -- Recorte dos coeficientes originais M01/M02; S01 nao e apresentado para este destino.
\begin{frame}{OLS: conversão para pastagem}
\small Coeficientes (erros-padrão), em pp. Preços e aptidões em desvios-padrão.
\par\smallskip
\begingroup
\fontsize{9.5}{11.4}\selectfont
\renewcommand{\arraystretch}{1.00}
\setlength{\tabcolsep}{4pt}
\begin{tabularx}{\textwidth}{L c c}
\toprule
 & \textbf{M01: múltipla} & \textbf{M02: fiscalização}\\
\midrule
Preço da soja & \est{-0{,}4809}{0,2845} & \est{-0{,}2063}{0,2507}\\
Preço do milho & \est{0{,}4308}{0,2710} & \est{0{,}1934}{0,2340}\\
Aptidão da soja & \est{0{,}1773}{0,0722} & \est{0{,}1179}{0,0613}\\
Aptidão do milho & \est{0{,}1527}{0,1001} & \est{0{,}2105}{0,1273}\\
Autos de flora / 10 & -- & \est{-0{,}0355}{0,0279}\\
\midrule
Controles de bioma & Sim & Sim\\
Observações & 2.772 & 1.140\\
Municípios / anos & 137 / 24 & 128 / 23\\
$R^2$ & 0,1272 & 0,0681\\
\bottomrule
\end{tabularx}
\endgroup
\par\smallskip
\footnotesize Modelos originais, com constante, bioma e controle adicional de preço. Apenas parte dos coeficientes é exibida.
\par\smallskip\small
M01: aptidão da soja positiva ($p=0{,}022$).
Em M02, os intervalos dos coeficientes exibidos incluem zero.
% \fonte{testes\_chap1, modelos\_coeficientes.csv e modelos\_resumo.csv.}
\end{frame}
%% COMENTÁRIO:
% (1) Agora eu faço a mesma leitura para a conversão de vegetação nativa em pastagem. Mantenho as demais variáveis constantes, com preços em dólares de 2024 e diferenças de um desvio-padrão, como no slide anterior.
% (2) No M01, para cada US$ 132,44 por tonelada a mais no preço da soja do ano anterior, a taxa estimada de conversão para pastagem é 0,4809 ponto percentual menor. Para cada US$ 73,57 por tonelada a mais no preço do milho do ano anterior, a taxa estimada é 0,4308 ponto percentual maior.
% (3) Entre municípios, uma aptidão para soja 379,58 kg por hectare maior está associada a uma taxa estimada 0,1773 ponto percentual maior. Para o milho, uma aptidão 778,52 kg por hectare maior está associada a uma taxa 0,1527 ponto percentual maior. São rendimentos potenciais agrícolas em peso seco dos grãos, não medidas da produtividade da pastagem.
% (4) Num exemplo hipotético de M01, se a taxa prevista fosse de 1%, essa diferença de aptidão para soja corresponderia a uma taxa de aproximadamente 1,1773%. A aptidão da soja é o único coeficiente exibido em M01 com p abaixo de 5%: p igual a 0,022, sem correção por múltiplos testes. Há evidência de associação positiva nesse modelo, mas isso não identifica uma relação causal.
% (5) No M02, a mesma diferença de aptidão para soja está associada a uma taxa 0,1179 ponto percentual maior. O p-valor passa para 0,067, e o intervalo de confiança de 95% inclui zero. Isso também acontece com os intervalos dos demais coeficientes exibidos em M02, mostrando a incerteza dessas estimativas.
% (6) Para a fiscalização, cada dez autos de flora adicionais no ano anterior estão associados a uma taxa estimada de conversão para pastagem 0,0355 ponto percentual menor. Essa associação é incerta e não permite concluir se a fiscalização foi eficaz.
% (7) O R quadrado passa de 12,72% em M01 para 6,81% em M02, mas a amostra também cai de 2.772 para 1.140 observações. Por isso, eu não atribuo a mudança dos coeficientes ou do ajuste apenas à inclusão da fiscalização: também mudou o conjunto de observações usado na regressão.

% ================================================
% #8 -- Avaliacao solicitada: utilidade associativa e insuficiencia causal.
\begin{frame}{A OLS serve para esta análise?}
\begin{block}{Serve como referência descritiva e associativa}
Permite medir associações em pontos percentuais, comparar ajustes e mostrar a sensibilidade à amostra. O baixo $R^2$, por si só, não invalida o estimador.
\end{block}
\medskip
\textbf{As três OLS não bastam para testar causalmente H1.}
\small
\begin{itemize}
\setlength{\itemsep}{2mm}
\item A interpretação causal exige justificar a exogeneidade dos regressores $X$: $E[u^d_{it}\mid X]=0$. As defasagens, sozinhas, não a asseguram.
\item Omissões municipais e temporais, autuações endógenas e medidas incompletas de retorno comprometem essa interpretação.
\item Erros-padrão agrupados tratam parte da dependência dos erros, mas não corrigem endogeneidade. Poucos anos limitam a inferência.
\end{itemize}
\medskip
Não há diagnóstico documentado de resíduos que permita declarar uma reprovação geral da OLS. A limitação central é o desenho e o alcance das medidas.
% \fonte{testes\_chap1/LEIA\_ME.md; testes7/CONCLUSAO\_GERAL.md; \href{https://cameron.econ.ucdavis.edu/research/Cameron_Miller_JHR_2015_February.pdf}{Cameron e Miller (2015)}.}
\end{frame}
%% COMENTÁRIO:
% (1) A resposta é que a OLS serve como ponto de partida. Ela me ajuda a medir as associações e a ver como os resultados mudam quando acrescento variáveis ou altero a amostra. Um R quadrado baixo, sozinho, não torna essas estimativas inválidas.
% (2) Para interpretar um coeficiente como efeito causal, eu preciso justificar a condição escrita no slide. Em termos simples, os fatores que ficaram de fora e afetam a conversão não podem estar relacionados às variáveis explicativas de modo a distorcer a comparação.
% (3) A fiscalização mostra bem o problema. Pode haver mais autos justamente onde houve mais infrações. Então, a quantidade de autos também pode estar respondendo ao processo que eu quero explicar. Usar a informação do ano anterior não elimina essa relação automaticamente.
% (4) Os erros-padrão agrupados levam em conta parte da dependência entre as observações, mas não corrigem esses problemas de causalidade. E, nestas OLS, eu tenho só 23 ou 24 anos para a inferência, o que também exige cautela.
% (5) Então, o limite principal está no que foi medido e na forma como a comparação foi construída. Preços não são retorno líquido e contar autos não mede diretamente a efetividade institucional. É por isso que eu uso essas regressões como evidência associativa.

% ================================================
% #9
\begin{frame}{Conclusões preliminares}
\begin{enumerate}
\setlength{\itemsep}{4mm}
\item As OLS mostram \textbf{associações brutas ou condicionais}, com sinais e precisão sensíveis à composição amostral. O sinal negativo dos preços não refuta H1.
\item Os modelos com efeitos fixos e interações oferecem evidência parcial de \textbf{heterogeneidade da resposta aos incentivos}, sobretudo na pastagem, mas permanecem sensíveis aos controles territoriais.
\item A hipótese completa \textbf{permanece em aberto}. Testá-la requer medidas defensáveis de retorno e efetividade institucional, além de uma estratégia de identificação compatível com as interações propostas.
\end{enumerate}
\medskip
\centering\color{verde}\textbf{OLS é um ponto de partida útil. Estes resultados não identificam efeitos causais.}
% \fonte{Síntese de testes\_chap1 e das rodadas testes2--testes7.}
\end{frame}
%% COMENTÁRIO:
% (1) Para fechar, eu vejo essas regressões como um diagnóstico inicial. Várias associações são imprecisas, e a comparação entre M01 e M02 mostra que preciso olhar com cuidado para a composição da amostra.
% (2) O sinal negativo do preço da soja é um resultado que eu preciso investigar. Ele não basta para refutar a hipótese, porque estas regressões não isolam o efeito causal do retorno esperado nem testam todas as interações previstas em H1.
% (3) Além destas tabelas, os testes complementares com efeitos fixos e interações sugerem que a associação entre preços e conversão varia com a aptidão e com o território, principalmente para pastagem. Mas esses resultados também mudam quando considero choques específicos de cada bioma ao longo do tempo.
% (4) Então, a hipótese completa continua em aberto. O próximo passo é melhorar as medidas de retorno esperado e de efetividade institucional e construir uma comparação que permita separar seus efeitos dos demais fatores. A OLS ajuda a organizar essa investigação, mas os resultados apresentados ainda não identificam efeitos causais.

\end{document}
